\begin{figure}[ht]
\centering
\begin{tikzpicture}[x=0.85cm, y=0.85cm, font=\small]

% Constraints: x1 <= 4, 2 x2 <= 12 (x2 <= 6), 3 x1 + 2 x2 <= 18.
% Feasible polygon vertices: (0,0) (4,0) (4,3) (2,6) (0,6).
% Plot region x1 in [0,6], x2 in [0,7].

% Axes
\draw[->, thick] (-0.2, 0) -- (6.6, 0) node[right, font=\scriptsize] {$x_{1}$};
\draw[->, thick] (0, -0.2) -- (0, 7.2) node[above, font=\scriptsize] {$x_{2}$};

% Tick labels
\foreach \t in {2, 4, 6} {
  \draw (\t, 0.08) -- (\t, -0.08) node[below, font=\scriptsize] {\t};
  \draw (0.08, \t) -- (-0.08, \t) node[left, font=\scriptsize] {\t};
}

% Feasible polygon (shaded)
\fill[engblue!12]
  (0,0) -- (4,0) -- (4,3) -- (2,6) -- (0,6) -- cycle;

% Constraint lines
% x1 = 4 (machine A)
\draw[engblack, thin] (4, -0.1) -- (4, 7.0);
\node[engblack, anchor=south, font=\scriptsize, rotate=90] at (4, 6.4)
  {\;$x_{1}=4$ (A)};
% x2 = 6 (machine B)
\draw[engblack, thin] (-0.1, 6) -- (6.4, 6);
\node[engblack, anchor=east, font=\scriptsize] at (6.3, 6.3)
  {$x_{2}=6$ (B)};
% 3 x1 + 2 x2 = 18 (machine C): passes (0,9) and (6,0); in window (2,6)-(6,0)
\draw[engblack, thin] (0, 9) -- (6, 0);
\node[engblack, anchor=west, font=\scriptsize] at (4.6, 2.3)
  {$3x_{1}+2x_{2}=18$ (C)};

% Objective level line through optimum: 3 x1 + 5 x2 = 36 -> x2 = (36 - 3 x1)/5
% at x1=2 -> x2=6; at x1=6 -> x2=3.6; at x1=0 -> 7.2
\draw[engred, thick, dashed] (0, 7.2) -- (6, 3.6);
\node[engred, anchor=west, font=\scriptsize] at (0.2, 6.95)
  {$3x_{1}+5x_{2}=36$};

% Vertices
\foreach \p/\lab/\pos in {%
  {(0,0)}/{}/below left,
  {(4,0)}/{}/below,
  {(4,3)}/{}/right,
  {(0,6)}/{}/left} {
  \fill[engblack] \p circle (1.6pt);
}
% Optimal vertex
\fill[engred] (2,6) circle (2.4pt);
\node[engred, anchor=south west, font=\scriptsize] at (2.1, 6.05)
  {$\vect{x}^{*}=(2,6)$, $z=36$};

\end{tikzpicture}
\caption[Graphical solution of a two-variable LP]{The feasible
polygon of the two-variable production LP, the intersection of the
three resource half-spaces with the first quadrant. The objective
$3x_{1}+5x_{2}$ is constant on the dashed parallel lines; pushing
the line outward until it last touches the polygon locates the
optimum at the vertex $\vect{x}^{*}=(2,6)$, where the machine-B
and machine-C constraints are both active and the machine-A
constraint is slack. The simplex method reaches the same vertex by
walking the edges $(0,0)\to(0,6)\to(2,6)$, each step increasing the
objective. In $n$ dimensions the polygon becomes a polyhedron with
exponentially many vertices, but the optimum is still attained at a
vertex whenever one exists.}
\label{fig:vol02:ch15:lp-graphical}
\end{figure}
